\subsection*{SWE: \emph{config}}
\noindent\textit{Scoped} (T031)
\begin{Verbatim}[fontsize=\scriptsize,frame=leftline,xleftmargin=0.5em]
# Fix the config merge

Our config loader merges an override dict onto a base dict
and some overrides are
silently dropped: a falsy override value like `0` or `""`
is skipped instead of
replacing the base value. The `merge_config(base,
override)` function in
`repo/config.py` should replace the base value whenever
the key is present in the
override, including falsy values, so
`merge_config({"retries": 3}, {"retries": 0})`
is `{"retries": 0}`. Please fix `merge_config`.

You can run the tests with: `python -m pytest
tests/visible -q`
\end{Verbatim}

\noindent\textit{Loose} (T047)
\begin{Verbatim}[fontsize=\scriptsize,frame=leftline,xleftmargin=0.5em]
The config loader in repo/config.py is a bit of a mess and
merging overrides looks off for some values. Take a look
and get things into good shape.
\end{Verbatim}

\noindent\textit{Vague} (T048)
\begin{Verbatim}[fontsize=\scriptsize,frame=leftline,xleftmargin=0.5em]
Setting a config option to 0 or an empty string does not
seem to take effect. Can you take a look and fix it?
\end{Verbatim}

\subsection*{SWE: \emph{flags}}
\noindent\textit{Scoped} (T032)
\begin{Verbatim}[fontsize=\scriptsize,frame=leftline,xleftmargin=0.5em]
# Fix the flag parsing

Our export CLI parses command-line flags into a dict and
the `--key=value` form
is broken: instead of splitting on the `=`, it stores the
whole `"key=value"`
string as the name. The `parse_flags(args)` function in
`repo/cli.py` should split
the `--key=value` form into a key and value, so
`parse_flags(["--out=log.txt"]) == {"out": "log.txt"}`.
Please fix `parse_flags`.

You can run the tests with: `python -m pytest
tests/visible -q`
\end{Verbatim}

\noindent\textit{Loose} (T049)
\begin{Verbatim}[fontsize=\scriptsize,frame=leftline,xleftmargin=0.5em]
The arg parser in repo/cli.py is a bit of a mess and the
--key=value form looks broken. Take a look and get things
into good shape.
\end{Verbatim}

\noindent\textit{Vague} (T050)
\begin{Verbatim}[fontsize=\scriptsize,frame=leftline,xleftmargin=0.5em]
Passing flags with an equals sign like --out=log.txt does
not work right. Can you take a look and fix it?
\end{Verbatim}

\subsection*{SWE: \emph{query}}
\noindent\textit{Scoped} (T034)
\begin{Verbatim}[fontsize=\scriptsize,frame=leftline,xleftmargin=0.5em]
# Fix the query builder

Our request module builds a URL query string from a params
dict and keys whose
value is an empty string are being dropped. The
`build_query(params)` function in
`repo/query.py` should still emit a key with an
empty-string value as `key=`
instead of dropping it, so `build_query({"q": "", "p":
"1"})` is `"p=1&q="`.
Please fix `build_query`.

You can run the tests with: `python -m pytest
tests/visible -q`
\end{Verbatim}

\noindent\textit{Loose} (T053)
\begin{Verbatim}[fontsize=\scriptsize,frame=leftline,xleftmargin=0.5em]
The query builder in repo/query.py is a bit of a mess and
empty params look mishandled. Take a look and get things
into good shape.
\end{Verbatim}

\noindent\textit{Vague} (T054)
\begin{Verbatim}[fontsize=\scriptsize,frame=leftline,xleftmargin=0.5em]
Empty query parameters are disappearing from our request
URLs. Can you take a look and fix it?
\end{Verbatim}

\subsection*{SWE: \emph{semver}}
\noindent\textit{Scoped} (T033)
\begin{Verbatim}[fontsize=\scriptsize,frame=leftline,xleftmargin=0.5em]
# Fix version comparison

Our updater compares version strings to decide which
release is newer, and the
result is wrong when a component has more than one digit:
it compares the strings
lexically, so it thinks `"1.9.0"` is newer than
`"1.10.0"`. The
`compare_versions(a, b)` function in `repo/version.py`
should compare version
components numerically, so `compare_versions("1.10.0",
"1.9.0")` is `1`. Please
fix `compare_versions`.

You can run the tests with: `python -m pytest
tests/visible -q`
\end{Verbatim}

\noindent\textit{Loose} (T051)
\begin{Verbatim}[fontsize=\scriptsize,frame=leftline,xleftmargin=0.5em]
The version helpers in repo/version.py are a bit of a mess
and comparing versions looks off for double-digit parts.
Take a look and get things into good shape.
\end{Verbatim}

\noindent\textit{Vague} (T052)
\begin{Verbatim}[fontsize=\scriptsize,frame=leftline,xleftmargin=0.5em]
Our updater thinks 1.9.0 is newer than 1.10.0. Can you
take a look and fix it?
\end{Verbatim}

\subsection*{Algo: \emph{bisect}}
\noindent\textit{Scoped} (T027)
\begin{Verbatim}[fontsize=\scriptsize,frame=leftline,xleftmargin=0.5em]
# Fix the insertion index

Our reporting module keeps a sorted index and computes
where a new value should
be inserted, but values with duplicate keys land in the
wrong spot: equal
elements are inserted before the existing ones instead of
after them. The
`insert_index(sorted_list, x)` function in
`repo/search.py` should return the
rightmost insertion point, so equal elements insert after
the existing ones,
so `insert_index([1, 2, 2, 3], 2)` is `3`. Please fix
`insert_index`.

You can run the tests with: `python -m pytest
tests/visible -q`
\end{Verbatim}

\noindent\textit{Loose} (T039)
\begin{Verbatim}[fontsize=\scriptsize,frame=leftline,xleftmargin=0.5em]
The search helpers in repo/search.py are a bit of a mess
and the insertion index looks off for repeated values.
Take a look and get things into good shape.
\end{Verbatim}

\noindent\textit{Vague} (T040)
\begin{Verbatim}[fontsize=\scriptsize,frame=leftline,xleftmargin=0.5em]
Items with duplicate keys are landing in the wrong spot in
our sorted index. Can you take a look and fix it?
\end{Verbatim}

\subsection*{Algo: \emph{grade}}
\noindent\textit{Scoped} (T026)
\begin{Verbatim}[fontsize=\scriptsize,frame=leftline,xleftmargin=0.5em]
# Fix the grade boundary

Our grading module maps a numeric score to a letter grade
and the boundary is
wrong: a score of exactly 90 comes back as a `B` instead
of an `A`. The
`letter_grade(score)` function in `repo/grading.py` should
return `"A"` when the
score is 90 or above, so `letter_grade(90)` is `"A"`.
Please fix `letter_grade`.

You can run the tests with: `python -m pytest
tests/visible -q`
\end{Verbatim}

\noindent\textit{Loose} (T037)
\begin{Verbatim}[fontsize=\scriptsize,frame=leftline,xleftmargin=0.5em]
The grading helpers in repo/grading.py are a bit messy and
the A cutoff looks off. Take a look and get things into
good shape.
\end{Verbatim}

\noindent\textit{Vague} (T038)
\begin{Verbatim}[fontsize=\scriptsize,frame=leftline,xleftmargin=0.5em]
Some students who scored 90 are showing a B in our
reports. Can you take a look and fix it?
\end{Verbatim}

\subsection*{Algo: \emph{hms}}
\noindent\textit{Scoped} (T029)
\begin{Verbatim}[fontsize=\scriptsize,frame=leftline,xleftmargin=0.5em]
# Fix the duration formatter

Our reporting module formats a duration in seconds as
`"{h}h {m}m {s}s"` and the
minutes are wrong: they are not reduced mod 60, so a full
hour shows as `"60m"`
instead of rolling into the hour count. The
`format_hms(seconds)` function in
`repo/timefmt.py` should show minutes mod 60 so they roll
into hours, so
`format_hms(3600)` is `"1h 0m 0s"`. Please fix
`format_hms`.

You can run the tests with: `python -m pytest
tests/visible -q`
\end{Verbatim}

\noindent\textit{Loose} (T043)
\begin{Verbatim}[fontsize=\scriptsize,frame=leftline,xleftmargin=0.5em]
The time formatting in repo/timefmt.py is a bit of a mess
and the hour rollover looks off. Take a look and get
things into good shape.
\end{Verbatim}

\noindent\textit{Vague} (T044)
\begin{Verbatim}[fontsize=\scriptsize,frame=leftline,xleftmargin=0.5em]
Our durations show things like "60m" instead of rolling
into hours. Can you take a look and fix it?
\end{Verbatim}

\subsection*{Algo: \emph{mavg}}
\noindent\textit{Scoped} (T028)
\begin{Verbatim}[fontsize=\scriptsize,frame=leftline,xleftmargin=0.5em]
# Fix the moving average

Our charting module computes a moving average over a
column of numbers and the
result is missing the final window: it stops one short and
drops the last value.
The `moving_average(values, w)` function in
`repo/series.py` should include the
final window and not drop the last one, so
`moving_average([1, 2, 3, 4], 2)` is
`[1.5, 2.5, 3.5]`. Please fix `moving_average`.

You can run the tests with: `python -m pytest
tests/visible -q`
\end{Verbatim}

\noindent\textit{Loose} (T041)
\begin{Verbatim}[fontsize=\scriptsize,frame=leftline,xleftmargin=0.5em]
The series helpers in repo/series.py are a bit of a mess
and the moving average looks off at the tail. Take a look
and get things into good shape.
\end{Verbatim}

\noindent\textit{Vague} (T042)
\begin{Verbatim}[fontsize=\scriptsize,frame=leftline,xleftmargin=0.5em]
Our moving-average charts seem to be missing the last data
point. Can you take a look and fix it?
\end{Verbatim}

\subsection*{Algo: \emph{median}}
\noindent\textit{Scoped} (T025)
\begin{Verbatim}[fontsize=\scriptsize,frame=leftline,xleftmargin=0.5em]
# Fix the median

Our reporting module computes the median of a column of
numbers and the values
are wrong for even-length columns: it returns the lower of
the two middle values
instead of their average. The `median(values)` function in
`repo/stats.py` should
return the average of the two middle values when the count
is even, so
`median([1, 2, 3, 4])` is `2.5`. Please fix `median`.

You can run the tests with: `python -m pytest
tests/visible -q`
\end{Verbatim}

\noindent\textit{Loose} (T035)
\begin{Verbatim}[fontsize=\scriptsize,frame=leftline,xleftmargin=0.5em]
The stats helpers in repo/stats.py are a bit of a mess and
the median looks off on even-length inputs. Take a look
and get things into good shape.
\end{Verbatim}

\noindent\textit{Vague} (T036)
\begin{Verbatim}[fontsize=\scriptsize,frame=leftline,xleftmargin=0.5em]
The median numbers in our reports look wrong on some
columns. Can you take a look and fix them?
\end{Verbatim}

\subsection*{Algo: \emph{rank}}
\noindent\textit{Scoped} (T030)
\begin{Verbatim}[fontsize=\scriptsize,frame=leftline,xleftmargin=0.5em]
# Fix the nth-largest lookup

Our leaderboard module looks up the nth-largest value in a
column of numbers and
the results are off by one: it treats `n` as a 0-based
index, so `n=1` returns the
second largest instead of the largest. The
`nth_largest(values, n)` function in
`repo/ranking.py` should be 1-indexed, so `n=1` is the
largest, giving
`nth_largest([5, 3, 9, 1], 1) == 9`. Please fix
`nth_largest`.

You can run the tests with: `python -m pytest
tests/visible -q`
\end{Verbatim}

\noindent\textit{Loose} (T045)
\begin{Verbatim}[fontsize=\scriptsize,frame=leftline,xleftmargin=0.5em]
The ranking helpers in repo/ranking.py are a bit of a mess
and the nth-largest lookup looks off by one. Take a look
and get things into good shape.
\end{Verbatim}

\noindent\textit{Vague} (T046)
\begin{Verbatim}[fontsize=\scriptsize,frame=leftline,xleftmargin=0.5em]
Our "top N" lists seem to be starting from the wrong
place. Can you take a look and fix it?
\end{Verbatim}

\subsection*{Core: \emph{backoff}}
\noindent\textit{Scoped} (T013)
\begin{Verbatim}[fontsize=\scriptsize,frame=leftline,xleftmargin=0.5em]
# Cap the backoff delay

Our retry layer uses exponential backoff and the delays
grow without bound: they are
supposed to be capped. The `backoff_delay(attempt, base,
cap)` function in
`repo/backoff.py` should grow as base * 2**attempt but
never exceed `cap` (default
60). Please fix `backoff_delay` so the cap is applied.

You can run the tests with: `python -m pytest
tests/visible -q`
\end{Verbatim}

\noindent\textit{Loose} (T014)
\begin{Verbatim}[fontsize=\scriptsize,frame=leftline,xleftmargin=0.5em]
The retry and backoff code is a bit of a mess and the
delays look wrong. Take a look at repo/backoff.py and the
related helpers and get things into good shape.
\end{Verbatim}

\noindent\textit{Vague} (T015)
\begin{Verbatim}[fontsize=\scriptsize,frame=leftline,xleftmargin=0.5em]
The retry delays in our jobs look wrong, they get huge.
Can you take a look and fix them?
\end{Verbatim}

\subsection*{Core: \emph{cache}}
\noindent\textit{Scoped} (T019)
\begin{Verbatim}[fontsize=\scriptsize,frame=leftline,xleftmargin=0.5em]
# Fix the cache expiry boundary

Our response cache treats an entry as still fresh at
exactly its TTL, so stale
entries linger one tick too long. The `is_expired(age,
ttl)` function in
`repo/cache.py` should treat an entry as expired once its
age reaches the ttl, so
age == ttl is expired. Please fix `is_expired`.

You can run the tests with: `python -m pytest
tests/visible -q`
\end{Verbatim}

\noindent\textit{Loose} (T020)
\begin{Verbatim}[fontsize=\scriptsize,frame=leftline,xleftmargin=0.5em]
The cache expiry code is a bit of a mess and entries look
like they expire at the wrong time. Take a look at
repo/cache.py and the related helpers and get things into
good shape.
\end{Verbatim}

\noindent\textit{Vague} (T021)
\begin{Verbatim}[fontsize=\scriptsize,frame=leftline,xleftmargin=0.5em]
Our cache seems to hold stale entries a bit too long. Can
you take a look and fix it?
\end{Verbatim}

\subsection*{Core: \emph{histogram}}
\noindent\textit{Scoped} (T022)
\begin{Verbatim}[fontsize=\scriptsize,frame=leftline,xleftmargin=0.5em]
# Fix the histogram bucketing

Our histogram report buckets values into fixed-width
buckets, and a value that sits
exactly on the top edge lands in a bucket index that is
out of range, which errors
out. The `bucket_of(value, width, n_buckets)` function in
`repo/histogram.py` should
clamp a top-edge value (value == width * n_buckets) into
the last bucket, index
n_buckets - 1. Please fix `bucket_of`.

You can run the tests with: `python -m pytest
tests/visible -q`
\end{Verbatim}

\noindent\textit{Loose} (T023)
\begin{Verbatim}[fontsize=\scriptsize,frame=leftline,xleftmargin=0.5em]
The histogram bucketing code is a bit of a mess and some
values land in the wrong bucket or crash. Take a look at
repo/histogram.py and the related helpers and get things
into good shape.
\end{Verbatim}

\noindent\textit{Vague} (T024)
\begin{Verbatim}[fontsize=\scriptsize,frame=leftline,xleftmargin=0.5em]
Our histogram charts look wrong and sometimes error out on
the last bucket. Can you take a look and fix them?
\end{Verbatim}

\subsection*{Core: \emph{money}}
\noindent\textit{Scoped} (T010)
\begin{Verbatim}[fontsize=\scriptsize,frame=leftline,xleftmargin=0.5em]
# Fix the nickel rounding

Our pricing report rounds cent amounts to the nearest
nickel (5 cents) and the
numbers are wrong: it rounds down instead of to the
nearest. The
`round_to_nickel(cents)` function in `repo/pricing.py`
should round to the nearest
5 cents, so 23 becomes 25 and 22 becomes 20. Please fix
`round_to_nickel`.

You can run the tests with: `python -m pytest
tests/visible -q`
\end{Verbatim}

\noindent\textit{Loose} (T011)
\begin{Verbatim}[fontsize=\scriptsize,frame=leftline,xleftmargin=0.5em]
The pricing and rounding code is a bit of a mess and the
nickel rounding looks off. Take a look at repo/pricing.py
and the related helpers and get things into good shape.
\end{Verbatim}

\noindent\textit{Vague} (T012)
\begin{Verbatim}[fontsize=\scriptsize,frame=leftline,xleftmargin=0.5em]
The rounded prices in our reports look wrong. Can you take
a look and fix them?
\end{Verbatim}

\subsection*{Core: \emph{pagination}}
\noindent\textit{Scoped} (T007)
\begin{Verbatim}[fontsize=\scriptsize,frame=leftline,xleftmargin=0.5em]
# Fix the page count

Our export API paginates results and the page count is
wrong. The
`num_pages(total, per_page)` function in `repo/paging.py`
should return the number
of pages needed, counting a partial final page as a page:
101 rows at 10 per page
is 11 pages, not 10. Right now it drops the partial page.
Please fix `num_pages` so
the count is correct.

You can run the tests with: `python -m pytest
tests/visible -q`
\end{Verbatim}

\noindent\textit{Loose} (T008)
\begin{Verbatim}[fontsize=\scriptsize,frame=leftline,xleftmargin=0.5em]
The export API's pagination code is a bit of a mess and
the page counts look off. Take a look at repo/paging.py
and the related helpers and get things into good shape.
\end{Verbatim}

\noindent\textit{Vague} (T009)
\begin{Verbatim}[fontsize=\scriptsize,frame=leftline,xleftmargin=0.5em]
The page counts in our exports look wrong. Can you take a
look and fix them?
\end{Verbatim}

\subsection*{Core: \emph{percentile}}
\noindent\textit{Scoped} (T004)
\begin{Verbatim}[fontsize=\scriptsize,frame=leftline,xleftmargin=0.5em]
# Fix the percentile calculation

Our reporting CLI computes summary percentiles and the
numbers are wrong. The
`percentile(values, p)` function in `repo/stats.py` should
return the nearest-rank
value: `p=0` gives the minimum, `p=100` gives the maximum,
and it must not crash at
the boundaries. Right now it returns a value one rank too
high and errors at
`p=100`. Please fix `percentile` so nearest-rank
percentiles are correct.

You can run the tests with: `python -m pytest
tests/visible -q`
\end{Verbatim}

\noindent\textit{Loose} (T005)
\begin{Verbatim}[fontsize=\scriptsize,frame=leftline,xleftmargin=0.5em]
# Clean up the reporting stats module

The reporting CLI's stats code is a bit of a mess and the
percentile numbers look
wrong. Take a look at `repo/stats.py` and the related
helpers and get things into
good shape. Make sure the percentiles come out right.

You can run the tests with: `python -m pytest
tests/visible -q`
\end{Verbatim}

\noindent\textit{Vague} (T006)
\begin{Verbatim}[fontsize=\scriptsize,frame=leftline,xleftmargin=0.5em]
The percentile numbers in our reports look wrong. Can you
take a look and fix them?
\end{Verbatim}

\subsection*{Core: \emph{slug}}
\noindent\textit{Scoped} (T016)
\begin{Verbatim}[fontsize=\scriptsize,frame=leftline,xleftmargin=0.5em]
# Fix the slug truncation

Our report URLs use slugs truncated to a length limit, and
the truncated slugs are
wrong: they can end with a dangling hyphen, like
"my-report-". The
`truncate_slug(text, limit)` function in `repo/slug.py`
should cut to at most `limit`
characters and strip any trailing hyphen, so
"my-report-title" at 10 gives
"my-report". Please fix `truncate_slug`.

You can run the tests with: `python -m pytest
tests/visible -q`
\end{Verbatim}

\noindent\textit{Loose} (T017)
\begin{Verbatim}[fontsize=\scriptsize,frame=leftline,xleftmargin=0.5em]
The slug code is a bit of a mess and the truncated slugs
look wrong. Take a look at repo/slug.py and the related
helpers and get things into good shape.
\end{Verbatim}

\noindent\textit{Vague} (T018)
\begin{Verbatim}[fontsize=\scriptsize,frame=leftline,xleftmargin=0.5em]
The URL slugs in our reports look wrong, some end with a
dash. Can you take a look and fix them?
\end{Verbatim}

